\documentclass[10pt,a4paper]{amsart}
\usepackage[margin=19mm]{geometry}
\usepackage{amsmath,amssymb,mathtools}
\usepackage[scr=rsfs,cal=euler]{mathalfa}
\usepackage{xfrac}
\usepackage[unicode,hidelinks,bookmarks=false]{hyperref}
\newcommand{\ie}{i.e.\ }
\newcommand{\cf}{cf.\ }
\newcommand{\resp}{respectively\ }
\newcommand{\Subsection}{\subsection}
\providecommand{\nobreakdash}{\nobreak\hbox{-}}
\setlength{\emergencystretch}{2em}
\multlinegap0pt
\usepackage{bm}
\usepackage{bbm}
\usepackage{dsfont}
\usepackage{xspace}
\usepackage{cancel}
\usepackage{mathtools}
\usepackage{amsthm}
\makeatletter
\@ifundefined{proposition}{\newtheorem{proposition}{Proposition}}{}
\makeatother
\newcommand{\PP}{\mathbb P}
\title[The Reye geometry inside the 64 lines of the Schur quartic]{The Reye geometry inside the 64 lines of the Schur quartic}
\author{Pawe\l{} Nurowski}
\date{Source: arXiv:2609.10751v2 (CC BY 4.0)}
\begin{document}
\maketitle

\noindent\emph{Source and licence.} The Reye geometry inside the 64 lines of the Schur quartic. Version arXiv:2609.10751v2 (CC BY 4.0), distributed under the Creative Commons Attribution 4.0 International licence, as recorded in the arXiv licence metadata for this exact version and shown on the abstract page https://arxiv.org/abs/2609.10751v2. Full rights notice: licenses/arxiv-2609.10751-CC-BY-4.0.txt.\par\smallskip

\noindent\textit{Editorial scope.} Counted unit: the Proposition whose source label is \texttt{prop:d4-reye} in the article (source0.tex lines 197--199, source label \texttt{prop:d4-reye}), delivered together with the article's own complete original proof of it (source0.tex lines 201--227, one \texttt{proof} environment of 27 lines). No mathematics is written, repaired or re-derived: statement and proof are the article's own text line for line. Required context retained uncounted: none. Every definition, display and label of those lines is retained unchanged. Omitted from this package and not redistributed: 1--196, 200--200, 228--1144. Nothing inside a retained excerpt is omitted or altered: every retained body is delivered exactly as the article's lines are written, so no omission receipt of any kind accompanies this package. Citations: the retained excerpts carry no citation command at all, so no bibliography entry of the article is transcribed and no reference is redistributed. Reference adaptation: none was needed and none was made; every reference inside the delivered unit resolves inside it, so no unresolved reference or citation is produced. Printed slips kept literal: no printed word, symbol, hypothesis or number of the counted statement or of its proof was altered. Layout and class: the article's own class is \texttt{article} with packages including fontenc, inputenc, lmodern, amsmath, amssymb, amsthm, mathtools, microtype, geometry, hyperref; the wrapper is the lane's pinned \texttt{amsart} editorial wrapper at 10pt on A4, which restates the theorem-like environments the retained text needs and adds the article's own macro definitions that the retained text needs (\texttt{\textbackslash PP}), with extra wrapper packages (bm, bbm, dsfont, xspace, cancel, mathtools). The delivered numbering is the unit's own. Assets: no retained span contains a figure environment, a graphics-inclusion command, a PDF-page inclusion, a table or a TikZ picture; no image file is copied into this package, the package declares no asset dependency and no image file is redistributed with it. Licence: version arXiv:2609.10751v2 of this article is distributed under the Creative Commons Attribution 4.0 International licence, as recorded in the arXiv licence metadata for this exact version and shown on the abstract page https://arxiv.org/abs/2609.10751v2. A full rights notice is delivered at licenses/arxiv-2609.10751-CC-BY-4.0.txt. Characters: every retained excerpt is pure ASCII, so the delivered engine, XeLaTeX with its default Latin Modern text font, reports no missing character.
\begin{proposition}\label{prop:d4-reye}
The projective configuration obtained from $\Phi(D_4)$ in this way is the cube realization of the Reye configuration described in the Introduction.
\end{proposition}

\begin{proof}
Use homogeneous coordinates $[X:Y:Z:W]$ on $\PP^3$ and consider the invertible linear map
\[
T(x_1,x_2,x_3,x_4)
=(x_1-x_3,\;x_1+x_3,\;x_2+x_4,\;x_2-x_4).
\tag{2.2}
\]
Four antipodal root pairs are sent to
\[
\begin{aligned}
[e_1-e_3]&\mapsto[1:0:0:0],&
[e_1+e_3]&\mapsto[0:1:0:0],\\
[e_2+e_4]&\mapsto[0:0:1:0],&
[e_2-e_4]&\mapsto[0:0:0:1].
\end{aligned}
\tag{2.3}
\]
The first three points are the three points at infinity of the coordinate directions in the affine chart $W\neq0$, and the fourth is the origin, which will be the centre of the cube.  The remaining eight antipodal root pairs are sent precisely to
\[
[\varepsilon_1:\varepsilon_2:\varepsilon_3:1],
\qquad \varepsilon_1,\varepsilon_2,\varepsilon_3\in\{\pm1\},
\tag{2.4}
\]
the eight vertices of the cube $[-1,1]^3$.

It remains only to interpret the $16$ lines.  Among the sixteen projective lines obtained from the zero-sum triples, four pass through the point $[0:0:0:1]$, which is the centre of the affine cube whose eight vertices are given in~(2.4).  If one of these four lines contains a vertex $[a:b:c:1]$, then, because it also passes through the centre, it contains the opposite vertex $[-a:-b:-c:1]$.  The eight vertices form four opposite pairs, so these four lines are precisely the four body diagonals of the cube.  Each of the three points at infinity in~(2.3) lies on four further zero-sum lines; under~$T$ these are exactly the four parallel cube edges in the corresponding coordinate direction.  Thus the remaining twelve lines are the twelve edges of the cube.  This is exactly the Reye incidence configuration.  \qedhere
\end{proof}

\end{document}
